\documentclass{uofa-eng-assignment}
\usepackage{caption}
\usepackage{hyperref}
\hypersetup{
    colorlinks = false,
    linkbordercolor = {white},
}
\captionsetup[table]{name=Tabla, labelsep=period}
\captionsetup[figure]{name=Figura, labelsep=period}

\newcommand{\hpl}{\hyperlink}
\newcommand{\hpt}{\hypertarget}
\newcommand*{\name}{Joan Sebastian Ramirez Sanchez  }
\newcommand*{\id}{202223948}
\newcommand*{\course}{Control óptimo aplicado a modelos biológicos}
\newcommand*{\assignment}{Tarea 1 de 3}
\newcommand*{\dated}{\today}
\newcommand{\tf}{\textbf}

\begin{document}

\maketitle
\textbf{Ejercicio 8.2} Resolver
\[
\max_{u} \int_{0}^{1} \bigl(x(t) - u(t)\bigr)\, dt
\]
sujeto a
\[
x'(t) = 2u(t)\bigl(1 - u(t)\bigr), \quad x(0) = 0,
\]
\[
0 \leq u(t) \leq 1.
\]
Como vimos en estas secciones, el problema de este ejercicio radica en la frontera
que tiene nuestro control, es decir, el hecho de que 
$u(t)$ debe estar entre $0$ y $1$. Para resolver este problema, primero encontramos la función Hamiltoniana de nuestro sistema, que es
\[H(x, u, \lambda) = x - u + 2\lambda u(1 - u).\]
Luego, encontramos la función de costo$\lambda(t)$, que es
\[\lambda'(t) = -\frac{\partial H}{\partial x} = -1 \implies \lambda(t) = -t + C.\]
 siguiendo las condiciones tenemos que 
\[
\frac{\partial H}{\partial u} = -1 + 2\lambda(1 - 2u) = 0.
\]
despejando $u$ tenemos que
\[
u = \frac{2\lambda - 1}{4\lambda}.
\]
Ahora

\end{document}
